% ===================== 第 6 章 =====================
\section{双图治理与科学事件账本}

\subsection{为什么必须是两张图}

过程图回答「科学家是怎样完成这次研究的」，证据图回答「我们目前对这个问题知道什么」。合并成一张图会同时引发五类问题：工具调用步骤与学术证据混杂；Agent 的中间猜想被误读为事实；压缩推理过程时误删证据；一篇论文里的多个发现无法精确表示；界面无法清晰呈现争议结构。因此两张图在存储、权限与折叠规则上全部隔离。

\begin{itemize}[leftmargin=1.6em, itemsep=3pt]
\item \textbf{过程图 Process Graph}：任务、工具调用、失败路线、质询与决策，由记忆层管理，允许普通 Flush/Fold/Recall，节点永不自动成为正式证据。
\item \textbf{证据图 Evidence Graph}：10 类正式节点、12 类正式边（速查表见附录 A），只存经过审计的知识，只有图投影器能写。
\end{itemize}

\subsection{科学事件账本：一切正式修改的唯一入口}

七名科学家不得直接写证据图，所有正式修改必须先写入 \code{packages/evidence/sql\_ledger.py} 的 \code{SqlEventLedger}。账本设计了四条性质。

\begin{enumerate}[leftmargin=1.6em, itemsep=3pt]
\item \textbf{幂等、可重放。}唯一约束 \code{uq\_event\_idempotency} 加全序 \code{sequence}（\code{coalesce(max(sequence),0)+1}）。同一批事件重放任意次，结果相同。
\item \textbf{冲突可见。}同一幂等键提交不同 payload 时抛 \code{EventConflict}，任何改写历史的尝试都无法静默通过。
\item \textbf{并发串行化。}PostgreSQL 事务级咨询锁（\code{pg\_advisory\_xact\_lock}）保证并发提交不乱序。
\item \textbf{过程事件不升级。}仅过程类事件允许 \code{on\_conflict="skip"}，正式事件不享受此通道。
\end{enumerate}

幂等键由阶段、席位与位置派生（\code{PhaseContext.key()}，含超长查询的哈希折叠——这是修复过真实生产事故的细节：把 \code{uuid4()} 放进 payload 的轮次会在重放时与自己的幂等键冲突，现有测试专门盯防这一回退）。

\subsection{图投影器：证据图的唯一写入者}

\code{packages/evidence/sql\_projector.py} 的 \code{SqlGraphProjector} 顺序消费账本事件，机制包括：节点事件类型白名单（过程事件直接拒绝）；\code{node\_id\_for} 保证节点跨重放 id 稳定；写入边时对悬空目标抛 \code{ProjectionError}；checkpoint 单调不回退。投影器与应用身份的数据库权限划分已在第 3.4 节给出，此处不再重复。

\subsection{双图双向联动}

\begin{figure}[ht]
\centering
\resizebox{\linewidth}{!}{%
\begin{tikzpicture}[node distance=3mm and 10mm]
  \node[nd, minimum width=34mm, minimum height=12mm] (proc) {\textbf{过程图}\newline\tiny 思考 · 检索 · 失败路线 · 质询};
  \node[nda, right=of proc, minimum width=30mm, minimum height=12mm] (led) {\textbf{事件账本 + 证据门}\newline\tiny 追加写 · 幂等 · 六阶段审计};
  \node[ndd, right=of led, minimum width=34mm, minimum height=12mm] (ev) {\textbf{证据图}\newline\tiny 主张 · 发现 · 来源 · 盲点};
  \draw[flow] (proc.north) to[bend left=20] node[lbl,above]{① 结构化事件送审} (led.north);
  \draw[flow] (led.north) to[bend left=20] node[lbl,above]{② 过审后投影落图} (ev.north);
  \draw[flow] (ev.south) to[bend left=20] node[lbl,below]{③ 缺口生成盲点与新任务} (led.south);
  \draw[flowg] (led.south) to[bend left=20] node[lbl,below]{④ 新任务回灌私人记忆} (proc.south);
\end{tikzpicture}}
\caption{双图联动：正向通道把过程提炼为证据，回流通道把证据缺口变成下一轮过程；两个方向都经过账本，没有旁路写入。}
\label{fig:dualgraph}
\end{figure}

\textbf{正向（过程 $\rightarrow$ 证据）}：科学家完成「思考 $\rightarrow$ 检索 $\rightarrow$ 阅读 $\rightarrow$ 抽取发现 $\rightarrow$ 提交动作」整段过程，过程图全程记录，但只有结构化事件进入审计通道。\textbf{回流（证据 $\rightarrow$ 过程）}：证据图发现「支持某主张的七项发现全部来自自报告数据」这类缺口，生成 \code{Blindspot} 与新的调查任务，写入相关席位的私人记忆，开启增量调查。两条路径都经过账本，不存在绕过审计的写入。

\subsection{被反驳的节点淡化，但永不消失}

投影器的 upsert 只改 \code{status}、不删行；辩证折叠要求正反双方字段全部非空，缺字段就拒绝折叠（\code{dialectical\_fold.py}）；异议证书挂在事件本体上，报告组装无法把它「总结没」。前端证据地图遵循同一条规则：被反驳节点淡化显示（dimmed, never hidden）。对应回归测试包括 \code{test\_dissent\_preservation.py} 与 \code{test\_dialectical\_fold.py}。
